\problem{2.10}{Second excited harmonic oscillator state}
Introduce $q=\sqrt{m\omega/\hbar}\,x$ and the normalized
$\psi_0(x)=(m\omega/\pi\hbar)^{1/4}\ee^{-q^2/2}$.
Using
$a^\dagger=(q-\partial_q)/\sqrt2$ and
$\psi_2=(a^\dagger)^2\psi_0/\sqrt{2!}$,
\begin{equation}
a^\dagger\psi_0=\sqrt2\,q\psi_0,\quad
(a^\dagger)^2\psi_0=(2q^2-1)\psi_0,\quad
\boxed{\psi_2(x)=\frac{2q^2-1}{\sqrt2}\psi_0(x)}.
\end{equation}
The first three are $\psi_0=\psi_0(x)$ (even, no
nodes), $\psi_1=\sqrt2q\psi_0$ (odd, one node $q=0$), and
$\psi_2$ (even, nodes at $q=\pm1/\sqrt2$). The inner products
$\avg{\psi_0|\psi_1}$ and $\avg{\psi_1|\psi_2}$ vanish
by oddness. The nontrivial pair is
\begin{align}
\int_{\mathbb R}\psi_0\psi_2\,\dd x
&=\frac1{\sqrt2\sqrt\pi}
\int_{-\infty}^\infty(2q^2-1)\ee^{-q^2}\dd q\\
&=\frac{2(\sqrt\pi/2)-\sqrt\pi}{\sqrt{2\pi}}=0.
\end{align}
\qcheck{Normalization of $\psi_2$ follows from
$\avg{(2q^2-1)^2}_{\psi_0}/2
=[4(3/4)-4(1/2)+1]/2=1$.}

\problem{2.11}{Lowest two oscillator states: direct expectations}
For $\psi_0\propto\ee^{-q^2/2}$ and
$\psi_1=\sqrt2q\psi_0$, the densities are even, so
$\avg x=\avg p=0$ (the functions are real and the
momentum integral is a boundary term).
Use the Gaussian moments
\begin{equation}
\frac1{\sqrt\pi}\int q^2\ee^{-q^2}\dd q=\frac12,\qquad
\frac1{\sqrt\pi}\int q^4\ee^{-q^2}\dd q=\frac34.
\end{equation}
For the ground state, $\avg{q^2}_0=1/2$.
For the first excited state,
$\avg{q^2}_1=2\avg{q^4}_0=3/2$. Since
$x=\ell q$, with $\ell=\sqrt{\hbar/(m\omega)}$,
\begin{equation}
\boxed{\avg{x^2}_0=\frac{\hbar}{2m\omega},\qquad
\avg{x^2}_1=\frac{3\hbar}{2m\omega}}.
\end{equation}
Acting with $-\hbar^2\partial_x^2$ on $\psi_0$ gives
$\hat p^2\psi_0=(m\hbar\omega-m^2\omega^2x^2)\psi_0$;
on $\psi_1$ gives
$\hat p^2\psi_1=(3m\hbar\omega-m^2\omega^2x^2)\psi_1$.
Thus
\begin{equation}
\boxed{\avg{p^2}_0=\frac{m\hbar\omega}{2},\qquad
\avg{p^2}_1=\frac{3m\hbar\omega}{2}}.
\end{equation}
The uncertainty products are respectively
$\hbar/2$ and $3\hbar/2$. The potential and
kinetic energies in these states satisfy
\begin{equation}
\avg V=\frac12m\omega^2\avg{x^2}=\frac12 E_n,\qquad
\avg T=\frac{\avg{p^2}}{2m}=\frac12 E_n,
\end{equation}
giving totals $E_0=\hbar\omega/2$ and
$E_1=3\hbar\omega/2$ as required.

\problem{2.12}{General oscillator expectations using ladder algebra}
Use
\begin{equation}
x=\sqrt{\frac{\hbar}{2m\omega}}(a+a^\dagger),\qquad
p=\ii\sqrt{\frac{m\hbar\omega}{2}}(a^\dagger-a),
\end{equation}
together with $\langle n|a|n\rangle=0$,
$a^\dagger a|n\rangle=n|n\rangle$ and
$aa^\dagger|n\rangle=(n+1)|n\rangle$.
Terms with $a^2$ or $(a^\dagger)^2$ change $n$ by two
and therefore vanish in diagonal matrix elements.
Consequently
\begin{align}
\avg x_n&=\avg p_n=0,\\
\boxed{\avg{x^2}_n}&=\boxed{\frac{\hbar}{m\omega}
\left(n+\frac12\right)},\\
\boxed{\avg{p^2}_n}&=\boxed{m\hbar\omega
\left(n+\frac12\right)},\\
\boxed{\avg T_n}&=\boxed{\frac12\hbar\omega
\left(n+\frac12\right)},\\
\boxed{\sigma_x\sigma_p}&=\boxed{\hbar\left(n+\frac12\right)}.
\end{align}
This obeys the uncertainty relation for all $n\ge0$,
with equality only in the ground state.

\problem{2.13}{Dynamics of a $3:4$ oscillator superposition}
Let $\Psi(x,0)=A[3\psi_0(x)+4\psi_1(x)]$.
Orthonormality gives $A=1/5$.
Writing $E_n=(n+1/2)\hbar\omega$,
\begin{equation}
\Psi(x,t)=\frac35\psi_0\ee^{-\ii\omega t/2}
+\frac45\psi_1\ee^{-3\ii\omega t/2}.
\end{equation}
Because the stationary functions are real,
\begin{equation}
|\Psi|^2=\frac9{25}\psi_0^2+\frac{16}{25}\psi_1^2+
\frac{24}{25}\psi_0\psi_1\cos(\omega t).
\end{equation}
Replacing $\psi_1$ by $\psi_2$ would produce a density
cross-term oscillating at $2\omega$ instead. Since
$\langle0|x|1\rangle=\sqrt{\hbar/(2m\omega)}$,
\begin{align}
\boxed{\avg x(t)}&=\boxed{\frac{24}{25}
\sqrt{\frac{\hbar}{2m\omega}}\cos(\omega t)},\\
\boxed{\avg p(t)}&=\boxed{-\frac{24}{25}
\sqrt{\frac{m\hbar\omega}{2}}\sin(\omega t)}.
\end{align}
Differentiating the second equation gives
$\dd\avg p/\dd t=-m\omega^2\avg x$,
which is Ehrenfest's theorem.
Energy outcomes are $E_0$ and $E_1$ with
probabilities $9/25$ and $16/25$, and
\begin{equation}
\avg H=\frac9{25}\frac12\hbar\omega+
\frac{16}{25}\frac32\hbar\omega
=\boxed{\frac{57}{50}\hbar\omega}.
\end{equation}

\problem{2.14}{Ground-state probability outside classical turning points}
The ground-state energy is $\hbar\omega/2$.
Classical turning points satisfy
$m\omega^2x_t^2/2=E_0$, giving $x_t=\ell$
with $\ell=\sqrt{\hbar/(m\omega)}$.
For $q=x/\ell$, the normalized ground-state density is
$\rho_q=\ee^{-q^2}/\sqrt\pi$. Thus
\begin{equation}
P(|x|>x_t)=\frac{2}{\sqrt\pi}
\int_1^\infty \ee^{-q^2}\dd q
=\operatorname{erfc}(1)
=\boxed{0.157\quad\text{(three significant figures)}}.
\end{equation}
The nonzero tail is a direct signature of quantum
penetration into the classically forbidden region.

\problem{2.15}{Fifth and sixth Hermite polynomials}
The Hermite equation $H_n''-2qH_n'+2nH_n=0$
admits the recursion
$H_{n+1}=2qH_n-2nH_{n-1}$,
with $H_0=1$, $H_1=2q$.
Starting from $H_3=8q^3-12q$ and
$H_4=16q^4-48q^2+12$,
\begin{align}
H_5&=2q(16q^4-48q^2+12)-8(8q^3-12q)\\
&=\boxed{32q^5-160q^3+120q},\\
H_6&=2q(32q^5-160q^3+120q)-10(16q^4-48q^2+12)\\
&=\boxed{64q^6-480q^4+720q^2-120}.
\end{align}
The leading coefficients are $2^5$ and $2^6$ respectively.

\problem{2.16}{Rodrigues, recursion, differentiation, generating function}
For physicists' Hermite polynomials,
\begin{equation}
H_n(q)=(-1)^n\ee^{q^2}\frac{\dd^n}{\dd q^n}
\ee^{-q^2},\quad
H_{n+1}=2qH_n-2nH_{n-1}.
\end{equation}
Repeated differentiation in Rodrigues' expression gives
\begin{equation}
H_3=8q^3-12q,\qquad
H_4=16q^4-48q^2+12.
\end{equation}
The recurrence then gives the $H_5$ and $H_6$
calculated in Problem 2.15. Direct differentiation checks
\begin{align}
H_5'&=160q^4-480q^2+120=10H_4,\\
H_6'&=384q^5-1920q^3+1440q=12H_5.
\end{align}
Finally expand the generating function
\begin{equation}
\ee^{-z^2+2qz}
=1+2qz+(2q^2-1)z^2+
\left(\frac43q^3-2q\right)z^3+O(z^4).
\end{equation}
Reading off the coefficient of $z^n/n!$ gives
\begin{equation}
\boxed{H_1=2q,\quad H_2=4q^2-2,\quad H_3=8q^3-12q}.
\end{equation}
The three identities are mutually consistent and provide
independent calculations of the same polynomial family.
